% ============================================================
% Chapter 5: Development
%   Section: Storage
% Included from chapters/05-development/development.tex
% Sources: docs/decisions/storage.md, graph.md, thresholds.md.
% ============================================================

\section{Storage, Journal and Per-Run Thresholds}
\label{sec:dev-storage}

\paragraph{The write path.} The analysis service writes the raw layer
after extraction, the processed layer after enrichment, and the
exploitation layer --- rewriting the processed record with the report
--- after verification (Section~\ref{sec:data-architecture}). Each
write is fail-soft on its own: a failure is reported as its own event
and does not stop the next layer, which only loses its parent link. A
storage error never discards a result that already cost a scrape, an
enrichment and a model call per claim. Writes go to a temporary file
and are moved into place atomically, because they run in the
background thread pool. The storage sits behind a five-method interface,
so moving the lake to a document database or object storage is one
class, not a pipeline change.

\paragraph{The journal.} The lake keeps results; what a run does on the
way --- the queries, every source returned, the rating each got --- used
to exist only in memory, so a crash lost exactly what was needed to
understand a bad verdict. Every job event is now appended to a file per
job as it happens. A job a restarted process no longer holds is still
answered from its file, and a file with no final line is reported as
interrupted rather than left looking as if it might still be running.

\paragraph{The analysis cache.} A repeat request for the same URL is
answered from a cache instead of re-running the pipeline. Two rules keep
it sound. Its key includes the run's \emph{effective} thresholds, not
only those that differ from the defaults: otherwise changing a default
would serve answers computed under the old value. And every entry is
stamped with a schema version that must change whenever the response
changes shape, because the cache never expires and never rewrites an
entry; the rule was broken in the week it was written, which is why it
is now a test that builds a real response and compares its keys with
the shape recorded for the current version.

\paragraph{Per-run thresholds.} Every tunable threshold lives in one
module, as a set of effective values for one run that default from the
environment, and a request may override any of them for that run only.
An override outside its range, or a misspelt one, is refused when the
request is made rather than when the job runs. Thresholds are passed to
each component per call, because the components are long-lived objects
shared by every request; the earlier practice of copying a setting into
a class attribute froze it at import time, so overrides had no effect,
and the tests that seemed to cover them reassigned the attribute
directly. A test now fails on any such copy. The lake records the
overrides of every run, so a stored verdict stays reproducible. The two
groups of scoring weights are the declared exception: they are set by
the environment only, because each group must sum to one and a caller
setting one weight alone would break that.

\paragraph{The graph sync.} The graph is written after the lake, in one
transaction per article, with short timeouts so that a stopped database
costs seconds rather than a minute per run. Hand-labelled facts reach
it only through a synchronisation that replays, in order, the source
catalogue, the labelled facts and the newest verified run per URL from
the lake. Every write merges on the schema's keys, so the sync can be
run any number of times. The labelling tool imports nothing from the
backend and does not write the graph itself.
